\begin{theoremEnd}[combined main,category=finite]{theorem}[Gaussian conditioning under a continuum sign restriction]
\CombinedResultCheck{thm:finite}{theorem}\label{thm:finite}
Under the finite Gaussian model in \cref{sec:finite-model}, let
$\boldsymbol a=(a_0,\ldots,a_m)^T$ denote the coefficient vector and fix
$r\in\{0,1,2\}$. Suppose an invertible linear change of coordinates
$\boldsymbol a\mapsto(\boldsymbol g,\gamma)\in\R^m\times\R$
gives $f_{\boldsymbol a}=f_{\boldsymbol g}+\gamma h_r$.
Here $h_r$ is a fixed function in the retained span, chosen to have a
uniformly positive $r$th derivative on $I$, and $\gamma$ is its scalar
coefficient. The curve $f_{\boldsymbol g}$ depends linearly on the
remaining coordinates $\boldsymbol g$. Assume these functions belong
to $C^r(I)$ with
$\min_{x\in I}\frac{\mathrm{d}^rh_r}{\mathrm{d}x^r}(x)>0$.
In these coordinates, the ordinary Gaussian posterior has the form
\begin{equation}\label{eq:conditional}
 \boldsymbol g\mid\boldsymbol y\sim N(\boldsymbol m_g,\boldsymbol V_g),
 \qquad \gamma\mid\boldsymbol g,\boldsymbol y\sim N(\mu(\boldsymbol g),s^2).
\end{equation}
Here $\boldsymbol V_g$ is positive definite, $\mu$ is affine, and
$s>0$ is constant in $\boldsymbol g$. The finite, attained threshold
\begin{equation}\label{eq:threshold}
 \tau_r(\boldsymbol g)=\max_{x\in I}
 \frac{-\frac{\mathrm{d}^rf_{\boldsymbol g}}{\mathrm{d}x^r}(x)}
      {\frac{\mathrm{d}^rh_r}{\mathrm{d}x^r}(x)}
\end{equation}
satisfies $f_{\boldsymbol a}\in\mathcal C_r$ exactly when
$\gamma\geq\tau_r(\boldsymbol g)$. This event has positive ordinary
posterior probability and a null boundary in coefficient space.
Its marginal density for $\boldsymbol g$ with
respect to Lebesgue measure on $\R^m$ is
\begin{equation}\label{eq:marginal}
 p(\boldsymbol g\mid\boldsymbol y,\mathcal C_r)=
 \frac{p(\boldsymbol g\mid\boldsymbol y)
 \Phi^{\mathrm c}((\tau_r(\boldsymbol g)-\mu(\boldsymbol g))/s)}
 {P(f_{\boldsymbol a}\in\mathcal C_r\mid\boldsymbol y)}.
\end{equation}
Given this draw of $\boldsymbol g$, draw $\gamma$ from the conditional normal in
\cref{eq:conditional} truncated below at $\tau_r(\boldsymbol g)$ and
apply the inverse coefficient map. Equivalently, for an independent
$U\sim\operatorname{Uniform}(0,1)$,
\begin{equation}\label{eq:reconstruction}
 \gamma=\mu(\boldsymbol g)-s\Phi^{-1}
 \bigl(U\Phi^{\mathrm c}((\tau_r(\boldsymbol g)-\mu(\boldsymbol g))/s)\bigr).
\end{equation}
This reconstructs the full conditioned coefficient law, under which
$\min_I\frac{\mathrm{d}^rf_{\boldsymbol a}}{\mathrm{d}x^r}>0$ almost surely.
\end{theoremEnd}

\begin{proofEnd}
\CombinedProofCheck{thm:finite}
Fix $r$, $I$, the invertible map, and the ordinary posterior under
the finite Gaussian model in \cref{sec:finite-model} and the
positive-direction hypotheses of
\cref{thm:finite}, and write
$b_r=\min_{x\in I}\frac{\mathrm{d}^rh_r}{\mathrm{d}x^r}(x)>0$.
The role of the positive direction is to turn every
derivative inequality into a lower bound on the same scalar coefficient:
division by $\frac{\mathrm{d}^rh_r}{\mathrm{d}x^r}$ gives the threshold equivalence,
and continuity on compact $I$, together with $b_r>0$, makes
\cref{eq:threshold} finite and attained. Since the invertible map
preserves positive definiteness of the posterior covariance, Gaussian conditioning
\citep[Section 2.2]{rasmussen2006gpml} gives
\cref{eq:conditional}, including a positive Schur complement $s^2$.

Integrating the restricted joint Gaussian density over the scalar
coefficient now leaves its conditional survival probability as a
weight on $\boldsymbol g$, which proves \cref{eq:marginal}.
This weight is strictly positive for every $\boldsymbol g$, so the
conditioning event has positive probability. Given a draw from the
weighted marginal, conditioning the Gaussian law of
$\gamma\mid\boldsymbol g,\boldsymbol y$ on
$\gamma\geq\tau_r(\boldsymbol g)$ gives the stated truncated normal,
whose survival-function inverse is \cref{eq:reconstruction}.
Applying the inverse coefficient map therefore recovers the original
coefficient law conditional on the sign restriction.

The minimum derivative is continuous in the coefficients, so coefficient
vectors with a positive minimum lie in the event's interior and those
with a negative minimum lie in its exterior. Conditional on
$\boldsymbol g$, the nondegenerate normal assigns zero
mass to equality at the threshold. At equality an attained maximizing
location has zero derivative, whereas above it
\[
 \min_I\frac{\mathrm{d}^rf_{\boldsymbol a}}{\mathrm{d}x^r}
 \geq(\gamma-\tau_r(\boldsymbol g))b_r>0.
\]
At equality, an arbitrarily small increase or decrease in $\gamma$
moves to either side of the event, identifying its coefficient boundary
as precisely the threshold graph.
Its zero Gaussian probability implies that the reconstructed law lies
strictly above the graph almost surely, proving the strict-derivative
assertion.
\end{proofEnd}
